\documentclass[paper=screen,mode=present,style=zysimple]{powerdot}

\usepackage{fancyvrb,verbatim,listings}
\usepackage{proof}
\usepackage[all]{xypic,xy}

\usepackage{latexsym,stmaryrd,amsmath,amsfonts}

%% \newtheorem{definition}{Definition}
%% \newtheorem{theorem}{Theorem}
%% \newtheorem{lemma}{Lemma}
\newtheorem{prop}{Proposition}
%% \newtheorem{corollary}{Corollary}
%% \newenvironment{proof}{{\sc Proof . }}{{\hfill $\square$}}
%% \newcommand{\myqed}{\hfill $\square$}
%% \newcommand{\ped}{\vspace{-1.6em}}
%% \newcommand{\note}[1]{\noindent\fbox{\parbox{\textwidth}{\em #1}}}

\newcommand{\abind}[2]{\mathtt{bind} \ (\, {#1} \,)  \ \mathtt{in} \ {#2}}
\newcommand{\acase}{\mathtt{case}}
\newcommand{\adv}{\mathcal{A}}
\newcommand{\aspecto}[1]{\xrightarrow{\; {#1} \;}}
\newcommand{\asplinear}{\mathfrak{l}}
\newcommand{\aspmodal}{\mathfrak{m}}
\newcommand{\aspnormal}{\mathfrak{n}}
%
\newcommand{\betareduce}{\leadsto_\beta}
\newcommand{\betareduceto}{\twoheadrightarrow_\beta}
\newcommand{\bit}{\mathtt{B}}
\newcommand{\bitone}{\mathtt{B}_1}
\newcommand{\bitzero}{\mathtt{B}_0}
%
\newcommand{\caseto}{\triangleright}
\newcommand{\cbind}[3]{\mathtt{bind} \ {#1} = {#2} \ \mathtt{in} \ {#3}}
\newcommand{\ccase}[4]{\mathtt{case} \; {#1} \; \mathtt{of} \; \cnil \caseto {#2}; \;
                                          \bitzero x_0 \caseto {#3}; \; \bitone x_1 \caseto {#4}}
\newcommand{\cequiv}{\simeq}
\newcommand{\cif}[3]{\mathtt{if} \; {#1} \; \mathtt{then} \; {#2} \; \mathtt{else} \; {#3}}
\newcommand{\ctrue}{\mathtt{true}}
\newcommand{\cfalse}{\mathtt{false}}
\newcommand{\cfix}{\mathtt{fix}\,}
\newcommand{\cfixin}[2]{\mathtt{fix}\; {#1}\; \mathtt{in} \; {#2}}
\newcommand{\cfold}{\mathtt{fold}\,}
\newcommand{\chd}{\mathtt{hd}}
\newcommand{\cinj}{\mathtt{inj}}
\newcommand{\clift}[1]{\lfloor{#1}\rfloor}
\newcommand{\clet}[3]{\mathtt{let} \ {#1} = {#2} \ \mathtt{in} \ {#3}}
\newcommand{\cletst}[3]{\mathtt{liftlet} \ {#1} = {#2} \ \mathtt{in} \ {#3}}
\newcommand{\cletrec}[3]{\mathtt{letrec} \ {#1} = {#2} \ \mathtt{in} \ {#3}}
\newcommand{\cnil}{\mathtt{nil}}
\newcommand{\conc}{\! \bullet \!}
\newcommand{\concat}{\pmb{\mathit{conc}}}
\newcommand{\context}{\mathbb{C}}
\newcommand{\cop}{\;\mathtt{op}\;}
\newcommand{\cpair}[1]{\langle{#1}\rangle}
\newcommand{\cproj}{\mathtt{proj}}
\newcommand{\crand}{\mathtt{rand}}
\newcommand{\crec}{\mathtt{rec}}
\newcommand{\croll}{\mathtt{roll}\,}
\newcommand{\crule}[1]{\mbox{{\textsl{#1}}}}
\newcommand{\cslrtype}{\Vdash}
\newcommand{\csucc}{\mathtt{S}}
\newcommand{\ctl}{\mathtt{tl}}
\newcommand{\cut}{\pmb{\mathit{cut}}}
\newcommand{\cval}{\mathtt{return}}
\newcommand{\cwhile}[2]{\mathtt{while}\; {#1} \; \mathtt{do} \; {#2} \; \mathtt{od}}
\newcommand{\cvalof}{\pmb{valueof}}
\newcommand{\cunfold}{\mathtt{unfold}\,}
\newcommand{\cunit}{\mathtt{unit}}
\newcommand{\cunroll}{\mathtt{unroll}\,}
\newcommand{\cupdate}{\pmb{update}}
\newcommand{\czero}{\mathtt{O}}
%
\newcommand{\deduces}{\vdash}
\newcommand{\defeq}{\stackrel{\mathrm{def}}{=}}
\newcommand{\delay}[1]{\left\lfloor{#1}\right\rfloor}
\newcommand{\distr}{\mathcal{D}}
%
\newcommand{\eagerreduce}{\xrightarrow{\mathrm{eager}}}
\newcommand{\Eagerreduce}{\stackrel{\mathrm{eager}}{\Reduce}}
\newcommand{\eqtest}{\stackrel{\mbox{\tiny ?}}{=}}
\newcommand{\Eval}[1]{{\left\llbracket{#1}\right\rrbracket}}
\newcommand{\evalcontext}{\mathbb{E}}
\newcommand{\Example}{\emgreen{\pmb{\textsc{Example}\;}}}
\newcommand{\expansion}[1]{\tilde{{#1}}}
%
\newcommand{\ford}{\mathrm{ord}}
%
\newcommand{\getval}{\leftarrow}
%
\newcommand{\head}{\pmb{\mathit{hd}}}
\newcommand{\termhole}{[\;]}
\newcommand{\hoaretriple}[3]{\{{#1}\} \; {#2} \; \{{#3}\}}
%
\newcommand{\imply}{\Rightarrow}
\newcommand{\indist}{\simeq}
\newcommand{\iteration}{\pmb{\mathit{it}}}
%
\newcommand{\lazyreduce}{\xrightarrow{\mathrm{lazy}}}
\newcommand{\Lazyreduce}{\lazyreduce\!\!\!\!\!\to}
\newcommand{\ldot}{\, . \, }
\newcommand{\linearto}{\multimap}
\newcommand{\logrel}{(\mathcal{R}_\tau)_{\mbox{$\tau$ type}}}
\newcommand{\lr}[1]{\mathcal{R}_{#1}}
\newcommand{\lub}{\sqcup}
%
\newcommand{\metalambda}{\underline{\lambda}}
\newcommand{\modal}{\square}
%
\newcommand{\nat}{\mathbb{N}}
\newcommand{\nReduce}{\nrightarrow\!\!\!\!\to}
\newcommand{\nreduceto}{\nrightsquigarrow}
\newcommand{\nullenv}{\varnothing}
%
\newcommand{\orc}{\mathcal{O}}
\newcommand{\oslrtype}{\Vdash}
%
\newcommand{\porder}{\sqsubseteq}
\newcommand{\prefix}{\pmb{\mathit{pref}}}
\newcommand{\proba}{\mathbf{Pr}}
%
\newcommand{\randlen}{\pmb{\mathit{rl}}}
\newcommand{\randnum}{\pmb{\mathit{rn}}}
\newcommand{\randstr}{\pmb{\mathit{rs}}}
\newcommand{\reduce}{\leadsto}
\newcommand{\reduceto}{\rightsquigarrow\hspace*{-6px}\shortrightarrow}
\newcommand{\Reduce}{\to\!\!\!\to}
\newcommand{\Reduceto}{\Downarrow}
%
\newcommand{\set}[1]{\{{#1}\}}
\newcommand{\setbits}{\mathbb{B}}
\newcommand{\slrtype}{\vdash}
\newcommand{\splt}{\pmb{\mathit{split}}}
\newcommand{\suffix}{\pmb{\mathit{suff}}}
\newcommand{\sysname}{\pmb{\mathcal{C}}}
%
\newcommand{\tail}{\pmb{\mathit{tl}}}
\newcommand{\tbits}{\mathsf{Bits}}
\newcommand{\tbool}{\mathsf{Bool}}
\newcommand{\tnat}{\mathsf{Nat}}
\newcommand{\tcomp}{\mathsf{T}}
\newcommand{\tlift}[1]{\lfloor{#1}\rfloor}
\newcommand{\tlist}{\mathsf{List}}
\newcommand{\tself}{\mathsf{self}}
\newcommand{\tstore}{\mathsf{Store}}
\newcommand{\tstream}{\mathsf{Stream}}
\newcommand{\tloc}{\mathsf{Loc}}
\newcommand{\ttcase}{\mathtt{case}}
\newcommand{\ttif}{\mathtt{if}}
\newcommand{\ttthen}{\mathtt{then}}
\newcommand{\ttelse}{\mathtt{else}}
\newcommand{\ttin}{\mathtt{in}}
\newcommand{\ttlet}{\mathtt{let}}
\newcommand{\ttof}{\mathtt{of}}
\newcommand{\tunit}{\mathsf{Unit}}
\newcommand{\Type}{\mathcal{T}\!\mathit{ype}}
%
\newcommand{\vbool}{\pmb{\mathbb{B}}}
\newcommand{\vff}{\mathfrak{f}\!\mathfrak{f}}
\newcommand{\vnat}{\pmb{\mathrm{N}}}
\newcommand{\vnil}{\epsilon}
\newcommand{\vtt}{\mathfrak{t}\!\mathfrak{t}}
%
\newcommand{\weakpre}{\pmb{wp}}




\pdsetup{lf={Theory and Practice of Functional Programming}}

\title{{\small \textsc{Theory and Practice of Functional Programming}}
\\[12pt]
{\Large \color{TitleColor} Simply Typed Lambda-Calculi (I)}
}
\author{\small \textbf{Dr. ZHANG Yu}
  \\[2pt] \small Institute of Software, Chinese Academy of Sciences 
  \\[5pt] \includegraphics[height=.04\slideheight]{iscas.eps} 
  \qquad \qquad 
  \includegraphics[height=.04\slideheight]{cas.eps}
}
\date{\scriptsize GUCAS, Beijing}

\newcommand{\RED}{\mathsf{RED}}

\begin{document}

\maketitle 

\begin{slide}{Adding type information}
\begin{itemize}
%% \item Almost all modern programming languages are typed.
%% \begin{itemize}
%% \item A strong type system can detect many errors \emblue{statically}.
%% \end{itemize}
\item Adding types to $\lambda$-calculus:
\[
\tau ::= b \mid \emred{\tau \to \tau}
\]
\begin{itemize}
\item The essential types are function types (also called arrow types). 
\item We often define some basic types for data like integers, booleans, etc. 
  $b$ ranges over a set of base types (type constants). 
\item Terms are normally the same as in untyped $\lambda$-calculus, but in some systems we also 
  write explicitly types for function arguments, i.e., $\lambda$-abstraction will be written as 
  $ \lambda x : \tau \ldot e$
\end{itemize}
\item \emblue{Order} of types:
\[
\begin{array}{l}
\ford(b) = 0 , \qquad \mbox{if $b$ is a base type} \\
%% \ford(\tau_1 \times \tau_2) = \max(\ford(\tau_1), \ford(\tau_2)) \\
%% \ford(\tau_1 + \tau_2) = \max(\ford(\tau_1), \ford(\tau_2)) \\
\ford(\tau_1 \to \tau_2) = \max(\ford(\tau_1)+1, \ford(\tau_2))
\end{array}
\]
\end{itemize}
\end{slide}

\begin{slide}{Type judgement}
\begin{itemize}
\item Type information of free variables are given by \emblue{typing contexts} and 
the type of a complete term can be justified.
\vspace*{-0.6em}
\[
x_1 : \tau_1,  \ldots, x_n : \tau_n \vdash e : \tau
\vspace*{-1.5em}
\]
\item \emblue{Typing rules} justify the derivation of typing assertions.
\[
\begin{prooftree}
x : \tau \in \Gamma
\justifies
\Gamma \vdash x : \tau 
\end{prooftree}
\qquad
\begin{prooftree}
\Gamma, x : \tau \vdash e : \tau' 
\justifies
\Gamma \vdash \lambda x: \tau \ldot e : \tau \to \tau'
\end{prooftree}
\qquad
\begin{prooftree}
\Gamma \vdash e : \tau' \to \tau
\quad 
\Gamma \vdash e' : \tau'
\justifies
\Gamma \vdash e\ e' : \tau
\end{prooftree}
\] 
\item Type safety
\begin{itemize}
\item \emblue{Subject reduction}: reductions do not change the type of expressions: 
For every expression $e$ such that $\Gamma \vdash e : \tau$ for some type $\tau$, 
if $e \reduce e'$, then $\Gamma \vdash e' : \tau$. 
\item \emblue{Progress}: all non-canonical closed terms can reduce. 
\end{itemize}
\end{itemize}
\end{slide}

\begin{slide}{G\"odel's system T}
\begin{itemize}
\item Types: 
\vspace*{-0.5em}
\[
\tau ::= \tnat \mid \tbool \mid \tau \to \tau 
\vspace*{-1.5em}
\]
\item Terms: 
\vspace*{-0.5em}
\[
e, e', \ldots ::= x \mid \lambda x \ldot e \mid e \ e' 
\mid \csucc e \mid \czero \mid \mathtt{T} \mid \mathtt{F} \mid \mathtt{D}\  e\ e'\ e''
\mid \mathtt{R}\ e\ e'\ e''
\vspace*{-1em}
\]
\begin{itemize}
\item $\mathtt O$ is the constant zero and $\mathtt S$ is the successor ($+1$) of type $\tnat \to \tnat$. 
\\[0.2em]
\item $\mathtt T, \mathtt F$ are boolean constants.
\\[0.2em]
\item $\mathtt D$ represents {\em if-then-else}:
\vspace*{-0.5em}
\[
\mathtt{D} \ \mathtt{T} \ e_1 e_2 \reduce e_1 , \qquad
\mathtt{D} \ \mathtt{F} \ e_1 e_2 \reduce e_2
\vspace*{-1.5em}
\]
\item $\mathtt R$ is the primitive recursor:
%  of type $\tnat \to \tau \to (\tnat \to \tau \to \tau) \to \tau$:
\vspace*{-0.5em}
\[
\mathtt{R} \ \mathtt{O} \ g \ f \reduce g , \qquad
\mathtt{R} \ (\mathtt{S}n) \ g \ f \reduce f \ n \ (\mathtt{R} \ n \ g \ f)
\]
\end{itemize}
\end{itemize}
\end{slide}

\begin{slide}{Programming in system T}
\begin{itemize}
\item Addition
$$\pmb{add} \defeq \lambda m \ldot \lambda n \ldot 
\mathtt{R} \ n \ m \ (\lambda x \ldot \lambda r \ldot \mathtt{S} \ r)
\vspace*{-2em}
$$ 
\item Iterator $\mathtt I$, which satisfies 
$$ \mathtt I \ \mathtt{O} \ g \ f \reduce g , \qquad
\mathtt{I} \ (\mathtt{S}n) \ g \ f \reduce f \ (\mathtt{R} \ n \ g \ f)
$$ 
Definition:
$$
\mathtt I \defeq \lambda n \ldot \lambda g \ldot \lambda f \ldot \mathtt{R} \ n \ g \ (\lambda x \ldot \lambda r \ldot f \ r)
\vspace*{-2em}
$$ 
\item The function class that is defined by system T is \emblue{primitive recursive functions}.
\begin{itemize}
\item Other PR functions include exponentiation and primality testing.
\item Using G\"odel numbering, we can operate on other objects like rational numbers.
\item Typical non-PR function: Ackermann function.
\end{itemize}
\end{itemize}
\end{slide}

\begin{slide}{Normal forms}
A term is {\em normal} if it does not contain any redex.  
$$\lambda x \ldot x, \quad \mathtt{S}, \quad \mathtt{S\ S\ O} \vspace*{-1em}$$
\begin{itemize}
\item \emblue{Head normal form (HNF)}: the out-most form is neither a redex nor an anstraction 
  with a reducible body.
\begin{itemize}
\item Variables, e.g., $x$;
\item Data values (constants), e.g., $\mathtt{O}, \mathtt{T}$;
\item Applications that do not reduce, e.g, 
  $x \ \mathtt{T} ,
  \csucc (\mathtt{R} \ \mathtt{O} \ \mathtt{O} \ (\lambda x \ldot \lambda r \ldot x))$
\item Abstractions with non-reducible bodies, e.g., 
  $\lambda x \ldot \mathtt{S}, \lambda x \ldot \lambda y \ldot x \ y$.
\end{itemize}
\item \emblue{Weak head normal form (WHNF)}: either a HNF or an abstration (no matter whether its body is reducible or not). E.g., 
  $\lambda x . \mathtt{R} \ \mathtt{O} \ \mathtt{T} \ x$.
\end{itemize}
\end{slide}

\begin{slide}{Canonical forms}
\begin{itemize}
\item Every closed normal term $e$ of system T must be in either of the following forms:
\begin{itemize}
\item $\mathtt{S}^n (\mathtt O)$ if $e$ has type $\tnat$ and $n = 0, 1, 2, \ldots$;
\item $\mathtt{T}$ or $\mathtt F$ if $e$ has type $\tbool$;
\item $\lambda x \ldot e'$, if $e$ has type $\tau_1 \to \tau_2$.
\end{itemize}
These are called \emblue{canonical forms} of system T. \\
The proof can be done by induction over structure of $e$.
\item Canonical forms of base types are often considered as {\em values} of computations.
\item All non-canonical closed terms can reduce.
\end{itemize}
\end{slide}


\begin{slide}{Normalization}
\begin{itemize}
% \item A term is {\em normal} if none of its subterm is of the form $(\lambda x: \tau \ldot e) e'$.
\item Normalization: Do normal forms always exist?
\begin{itemize}
\item \emblue{Weak normalization}: a term $e$ is {\em weakly normalizing} if there exists 
  a \emred{finite} sequence of reduction $e \to e_1 \to \cdots \to e_n$ and $e_n$ is normal. 
\\[0.2em]
\item \emblue{Strong normalization}: a term $e$ is {\em strongly normalizing} if there is 
  \emred{no infinite} sequence of reduction starting from $e$, i.e., 
  every reduction sequence starting from $e$ is finite (terminating).
\end{itemize}
\item System T is strongly normalizing. 
\end{itemize}
\end{slide}

\begin{slide}{Proof of weak normalization}
\emblue{Weak normalization}: for every term $e$, there exists a finite reduction sequence 
$e = e_0 \to e_1 \to \cdots \to e_n$ where $e_n$ is normal.
\vspace*{-0.5em}
\begin{itemize}
\item Observations
\vspace*{-0.5em}
\begin{itemize}
\item Simple induction over term structure does not work, since $\beta$-reduction often 
  breaks the structural induction. 
\item We need to find something that \emred{strictly decreases} with each single-step reduction.
\vspace*{-0.8em}
\end{itemize}
\item Degrees for terms and redexes
\vspace*{-0.5em}
\begin{itemize}
\item Degree of redexes: \vspace*{-1.2em}
\[
\partial ((\lambda x \ldot e) e') \defeq \partial (\tau' \to \tau)
\vspace*{-0.8em}
\]
where $\tau' \to \tau$ is the type of $\lambda x \ldot e$, and $\partial (\tau) $ is defined by
\vspace*{-0.8em}
\begin{eqnarray*}
\partial(b) & \defeq & 1, \quad \mbox{($b$ a base type)} 
\\
\partial(\tau \to \tau') & \defeq & \max(\partial(\tau), \partial(\tau')) + 1 
\\[-3.8em]
\end{eqnarray*}
\item Degree of terms: \vspace*{-0.8em}
\[
\delta(e) \defeq \max \set{\partial(r) \mid \mbox{$e$ is a redex and a subterm of $e$}}
\]
\end{itemize}
\end{itemize}
\end{slide}

\begin{slide}{Proof of weak normalization}
\begin{itemize}
\item A few lemmas 
\begin{itemize}
\item[1.] If $x$ is of type $\tau$, then $\delta(e[t /x]) \leq  \max(\delta(e), \delta(t), \delta(\tau))$.
\item[2.] If $e \to e'$, then $\delta(e') \leq \delta(e)$.
\item[3.] Let $r$ be a redex in $e$ with the maximal degree $n$, and supppose all of its redexes 
  have degrees that are strictly less than $n$. If $e'$ is the term obtained from $e$ 
  by converting the redex $r$, then $e'$ has strictly fewer redexes of degree $n$ than $e$.
\end{itemize}
\item For every term $e$, the proof makes a \emred{double induction over $(n, m)$}, where 
\[
\emblue{n = \delta(e), \qquad m = \mbox{the number of redexes of degree $n$ in $e$} .}
\]
The induction order between the pairs is defined by 
\[
(n_1, m_2) < (n_2, m_2) \mbox{ iff } n_1 < n_2 \mbox{ or } (n_1 = n_2 \;\&\; m_1 < m_2).
\]
\end{itemize}
\end{slide}

\begin{slide}{Proof of strong normalization}
\emblue{Strong normalization}: every reduction sequence of every term $e$ is finite.
\vspace*{-0.5em}
\begin{itemize}
\item Lemma: $e$ is strongly normalizing iff there exists a number $\nu(e)$ that bounds the length of 
  every reduction sequence from $e$.
\vspace*{-0.3em}
\item Define a \emred{type-indexed} set $\RED_\tau$ of terms by:
\vspace*{-0.5em}
\begin{itemize}
\item For a base type $b$, $e \in \RED_b$ if $e$ is strongly normalizing.
\item For arrow types, $e \in \RED_{\tau \to \sigma}$ iff for every $t \in \RED_\tau$, $e\ t \in \RED_\sigma$.
\vspace*{-0.5em}
\end{itemize}
\item A few properties about $\RED$:
\vspace*{-0.5em}
\begin{itemize}
\item[P1:] If $e \in \RED$, then $e$ is S.N.
\item[P2:] If $e \in \RED$ and $e \to e'$, then $e' \in \RED$.
\item[P3:] If all term $e'$ such that $e \to e'$ are in $\RED$ and $e$ is not an abstraction, 
  then $e \in \RED$.
\vspace*{-0.5em}
\end{itemize}
\item The rest of the proof consists of proving two lemmas:
\vspace*{-0.5em}
\begin{itemize}
\item If for every $t \in \RED$, $e[t / x] \in \RED$, then $\lambda x \ldot e \in \RED$.
\item For all $e, t \in \RED$, $e[t / x] \in \RED$.
\end{itemize}
\end{itemize}
\end{slide}


\end{document}
 

